\documentclass[aspectratio=169,11pt]{beamer}
\usepackage[utf8]{inputenc}
\usepackage[T1]{fontenc}
\usepackage[spanish,es-nodecimaldot]{babel}
\usepackage{amsmath,amssymb,mathtools}
\usepackage{booktabs,tabularx}
\usepackage{graphicx}
\usepackage{listings}
\usepackage{xcolor}
\usepackage{hyperref}

\usetheme{default}
\usecolortheme{seahorse}
\setbeamertemplate{navigation symbols}{}
\setbeamertemplate{footline}[frame number]
\setbeamerfont{frametitle}{size=\large}

\definecolor{qxblue}{HTML}{163A5F}
\definecolor{qxrust}{HTML}{A0442C}
\definecolor{qxgreen}{HTML}{276749}
\definecolor{qxgray}{HTML}{F1F4F7}
\setbeamercolor{structure}{fg=qxblue}
\setbeamercolor{alerted text}{fg=qxrust}

\lstdefinelanguage{Lean}{
  morekeywords={def,theorem,lemma,by,intro,have,apply,exact,rw,linarith,nlinarith,ring,norm_num},
  sensitive=true,
  morecomment=[l]{--},
  morecomment=[s]{/-}{-/}
}
\lstset{
  language=Lean,
  basicstyle=\ttfamily\scriptsize,
  keywordstyle=\color{qxblue}\bfseries,
  commentstyle=\color{qxgreen},
  columns=fullflexible,
  keepspaces=true,
  showstringspaces=false,
  frame=single,
  rulecolor=\color{black!20},
  backgroundcolor=\color{qxgray},
  breaklines=true,
  xleftmargin=0.4em,
  xrightmargin=0.4em
}

\newcommand{\sourcefoot}[1]{\vfill{\scriptsize\color{black!60}Fuente: #1}}
\newcommand{\handphoto}[2]{%
  \IfFileExists{#1}{%
    \includegraphics[width=\linewidth,height=.43\textheight,keepaspectratio]{#1}%
  }{%
    \fbox{\parbox[c][.36\textheight][c]{.90\linewidth}{\centering
      \textbf{Fotografía pendiente}\\[0.5em]
      \texttt{#1}\\[0.8em]
      #2}}%
  }%
}

\title{Quispe \& Xu (2026)}
\subtitle{\emph{Agentic Delegation and the Language Frontier of Software Developers}}
\author{Alejandro Ventura Ventura Meza}
\date{Repository 3 \textperiodcentered{} 4 de septiembre de 2026}

\begin{document}

%================= PORTADA ============================================
\begin{frame}[plain]
  \titlepage
  \begin{center}\small
    \href{https://github.com/alejandroventuraventurameza-tech/ai-03-quispe}
      {\texttt{github.com/alejandroventuraventurameza-tech/ai-03-quispe}}
  \end{center}
\end{frame}

%================= CITA ===============================================
\begin{frame}{1 \textperiodcentered{} La cita cambió el objeto de estudio}
  \begin{tabularx}{\textwidth}{@{}>{\bfseries}p{.20\textwidth}X@{}}
    \toprule
    Antes de verificar & ``Coding Beyond Your Training: Claude Code and the Technological Frontier of Software Developers'' \\
    \midrule
    Registro v2 & Alexander Quispe y Kevin Xu, \emph{Agentic Delegation and the Language Frontier of Software Developers} \\
    \midrule
    Fecha & arXiv registra v2 el 7 de julio. El PDF está fechado el 8 de julio de 2026 \\
    \midrule
    Muestra & 5,346 desarrolladores durante 28 meses \\
    \bottomrule
  \end{tabularx}

  \medskip
  El paper estudia una \alert{frontera de producción observada}. No demuestra
  que el desarrollador haya aprendido a programar solo en el nuevo lenguaje.

  \sourcefoot{Quispe y Xu (2026), arXiv:2605.25438v2, portada y resumen.}
\end{frame}

%================= PROBLEMA ===========================================
\begin{frame}{2 \textperiodcentered{} El problema del desarrollador}
  Para cada oportunidad desarrollador--lenguaje--mes, el desarrollador elige
  el modo con mayor excedente equivalente cierto:
  \[
    V^g_{ik,t}=\max_{m\in\mathcal M_g}V^m_{ik,t},
    \qquad
    Z^g_{ik,t}=\mathbf 1\!\left[V^g_{ik,t}\ge 0\right].
  \]

  \begin{columns}[T,onlytextwidth]
    \column{.48\textwidth}
    \textbf{Antes del agente}
    \[
      \mathcal M_1=\{S,C\}
    \]
    \begin{itemize}\small
      \item $S$: producción individual.
      \item $C$: conversación o autocompletado.
      \item La ganancia $\gamma s-r_C$ necesita habilidad previa.
    \end{itemize}

    \column{.48\textwidth}
    \textbf{Después de la adopción}
    \[
      \mathcal M_2=\{S,C,D\}
    \]
    \begin{itemize}\small
      \item $D$: ejecución delegada al agente.
      \item El valor depende de la habilidad general $a$ para especificar y verificar.
    \end{itemize}
  \end{columns}

  \sourcefoot{Sección 3, entorno y modos de producción.}
\end{frame}

%================= UMBRALES ===========================================
\begin{frame}{3 \textperiodcentered{} Umbrales de entrada}
  Un lenguaje se activa cuando la oportunidad $\omega$ supera el menor umbral
  disponible. Para un lenguaje desconocido, el supuesto de \emph{foothold}
  implica
  \[
    T^1=\min\{T^S,T^C\}=T^S.
  \]

  La delegación añade un umbral $T^D$. Su ventaja frente a producir solo es
  \begin{align*}
    B=T^S-T^D
    ={}&\lambda[az(A)-s\mu]-\kappa(a,s)-r_D \\
      &+\frac{\rho}{2}\left[
        \frac{(2\lambda-\lambda^2)s^2}{\pi}
        -\sigma_D^2(a,s,A)
      \right].
  \end{align*}

  \begin{center}
    \alert{$B>0$ equivale a $T^D<T^S$: delegar permite entrar antes.}
  \end{center}

  \sourcefoot{Sección 3, ecuaciones de umbrales y ventaja de delegación.}
\end{frame}

%================= BANDA ==============================================
\begin{frame}{4 \textperiodcentered{} La banda de activación}
  Si el lenguaje es desconocido y $B>0$, la Proposición 2 establece
  \[
    Z^2_{ik,t}-Z^1_{ik,t}
    =\mathbf 1\!\left[T^D_{ik,t}\le\omega_{ik,t}<T^S_{ik,t}\right].
  \]

  \begin{center}
    \begin{tabular}{c@{\hspace{1.2em}}c@{\hspace{1.2em}}c}
      $\omega<T^D$ & $T^D\le\omega<T^S$ & $T^S\le\omega$ \\
      \cmidrule(lr){1-1}\cmidrule(lr){2-2}\cmidrule(lr){3-3}
      Inactivo & \alert{Solo se activa con delegación} & Activo con ambos menús
    \end{tabular}
  \end{center}

  \medskip
  Con una CDF condicional continua $F$, la probabilidad de entrar en la banda es
  \[
    F(T^S_{ik,t})-F(T^D_{ik,t}).
  \]

  \small
  Condiciones que cargan el resultado: lenguaje inicialmente desconocido,
  $\gamma s-r_C\le0$, $B>0$ y continuidad de $F$ para la fórmula probabilística.

  \sourcefoot{Proposición 2 y ecuación de la banda de activación.}
\end{frame}

%================= EVIDENCIA ==========================================
\begin{frame}{5 \textperiodcentered{} Predicciones y evidencia de GitHub}
  \begin{tabularx}{\textwidth}{@{}X>{\raggedleft\arraybackslash}p{.25\textwidth}@{}}
    \toprule
    Resultado alrededor de la primera coautoría detectable con Claude Code & Estimación \\
    \midrule
    Lenguajes activos & $+2.5$ sobre una media previa de $0.9$ \\
    Lenguajes usados por primera vez & $+1.2$ \\
    Entropía de lenguajes & $+0.38$ \\
    Archivos reconstruidos & $57$ millones \\
    \bottomrule
  \end{tabularx}

  \medskip
  Los estudios de evento usan comparaciones con desarrolladores aún no tratados
  y varias pruebas de robustez. La adopción es voluntaria y puede coincidir con
  el inicio de un proyecto nuevo.

  \begin{center}
    \alert{La evidencia describe asociaciones en tiempo de evento, no un efecto causal definitivo.}
  \end{center}

  \sourcefoot{Resumen, Secciones 4--7 y tablas principales.}
\end{frame}

%================= IDENTIFICACION =====================================
\begin{frame}{6 \textperiodcentered{} Dos límites de identificación}
  \begin{columns}[T,onlytextwidth]
    \column{.48\textwidth}
    \textbf{La geometría no identifica al agente}

    Cualquier cambio que reduzca el umbral de entrada puede crear una banda:
    un subsidio, mejor documentación, un colaborador o un choque favorable de
    productividad.

    \medskip
    La Proposición 2 demuestra que la tecnología $D$ \emph{puede} producir la
    expansión. No demuestra que solo un agente pueda producirla.

    \column{.48\textwidth}
    \textbf{Selección en la adopción}

    Un desarrollador puede adoptar Claude Code precisamente cuando empieza un
    proyecto en Rust, Go u otro lenguaje nuevo.

    \medskip
    Los controles eliminan explicaciones mecánicas importantes. No observan el
    proyecto contrafactual que motivó la fecha de adopción.
  \end{columns}

  \sourcefoot{Auditoría propia del mecanismo e identificación, basada en las limitaciones del paper.}
\end{frame}

%================= COMPARACION ========================================
\begin{frame}{7 \textperiodcentered{} La decisión de modelado cambia la predicción}
  \begin{tabularx}{\textwidth}{@{}p{.22\textwidth}XX@{}}
    \toprule
      & Aouad--Lykouris--Zhong & Quispe--Xu \\
    \midrule
    Papel de la IA
      & Input sustituible dentro de $s+e+a$
      & Nuevo modo de ejecución delegada $D$ \\
    \addlinespace
    Margen central
      & La asistencia desplaza esfuerzo humano
      & La delegación reduce umbrales específicos por lenguaje \\
    \addlinespace
    Aprendizaje
      & Menor esfuerzo puede frenar acumulación de habilidad
      & Las interacciones generan señales que actualizan $\mu$ y $\pi$ \\
    \addlinespace
    Predicción
      & Menor esfuerzo y posible deskilling
      & Expansión de la frontera de producción \\
    \bottomrule
  \end{tabularx}

  \medskip
  No existe una conclusión universal de que la IA expanda o reduzca capacidad.
  La predicción depende de qué decisión humana sustituye, de qué opción nueva
  incorpora el modelo y de si la interacción genera aprendizaje.

  \sourcefoot{Comparación con el paper de la semana 1 solicitada en el issue del curso.}
\end{frame}

%================= PROP 3 / HAND 1 ====================================
\begin{frame}{8 \textperiodcentered{} Proposición 3: del flujo al stock}
  \begin{columns}[T,onlytextwidth]
    \column{.53\textwidth}
    Para un lenguaje inicialmente desconocido, el hazard constante $p^g_{ik}$
    produce
    \[
      \Pr(\mathcal T^g_{ik}>s)=(1-p^g_{ik})^{s+1}.
    \]
    Por complemento y suma sobre $k\in\mathcal U_i$,
    \[
      \Delta C_i(s)=\sum_{k\in\mathcal U_i}
      \left[(1-p^1_{ik})^{s+1}-(1-p^2_{ik})^{s+1}\right].
    \]
    Si $0\le p^1_{ik}\le p^2_{ik}\le1$, cada sumando es no negativo.

    \column{.43\textwidth}
    \centering
    \handphoto{hand/prop-3-page-1.jpg}{Supervivencia, complemento, suma y no negatividad.}
  \end{columns}

  \sourcefoot{Proposición 3 y Apéndice teórico.}
\end{frame}

%================= ENDPOINT / HAND 2 ==================================
\begin{frame}{9 \textperiodcentered{} El extremo $p^2=1$ rompe la estrictez}
  \begin{columns}[T,onlytextwidth]
    \column{.54\textwidth}
    Con frontera cerrada, $p^1=0$:
    \[
      d(s)=1-(1-p^2)^{s+1}.
    \]
    \[
      d(s+1)-d(s)=p^2(1-p^2)^{s+1},
    \]
    \[
      \Delta^2 d(s)=-(p^2)^2(1-p^2)^{s+1}.
    \]
    Ambas desigualdades son estrictas cuando $0<p^2<1$.

    En el extremo incluido por el texto,
    \[
      p^2=1 \quad\Longrightarrow\quad d(s)=1
      \quad\text{para todo }s\ge0.
    \]
    \alert{El efecto salta al inicio y luego permanece plano.}

    \column{.42\textwidth}
    \centering
    \handphoto{hand/prop-3-page-2.jpg}{Primera diferencia, concavidad y prueba del extremo.}
  \end{columns}
\end{frame}

%================= LEAN SPEC ==========================================
\begin{frame}[fragile]{10 \textperiodcentered{} Lean: de la ecuación al \texttt{Spec}}
  \textbf{Enunciado matemático corregido}
  \[
    p^1=0,\quad 0<p^2<1
    \quad\Longrightarrow\quad
    d(s)<d(s+1).
  \]

  \textbf{Traducción fuente en \texttt{PaperInterface.lean}}
\begin{lstlisting}
def cumulative_effect_strict_mono_closedFrontierSpec : Prop :=
  forall (p2 : Real) (s : Nat),
    0 < p2 -> p2 < 1 ->
    (1 - (1 - p2) ^ (s + 1)) <
    (1 - (1 - p2) ^ (s + 2))
\end{lstlisting}

  \small
  \texttt{Real} representa el hazard, \texttt{Nat} el horizonte discreto y las
  dos hipótesis abiertas excluyen los extremos donde la desigualdad deja de ser
  estricta. El \texttt{Spec} expone la afirmación que debe revisar una persona.
\end{frame}

%================= LEAN PROOF =========================================
\begin{frame}[fragile]{11 \textperiodcentered{} Lean: prueba, extremo y límite de la revisión}
  \begin{columns}[T,onlytextwidth]
    \column{.58\textwidth}
\begin{lstlisting}
theorem cumulative_effect_strict_mono_closedFrontier :
    cumulative_effect_strict_mono_closedFrontierSpec := by
  intro p2 s h0 h1
  have hb : 0 < 1 - p2 := by linarith
  have hpos : 0 < (1 - p2) ^ (s + 1) :=
    pow_pos hb (s + 1)
  rw [pow_succ]
  nlinarith [mul_pos hpos h0]
\end{lstlisting}

    \column{.38\textwidth}
    \small
    \textbf{Qué verifica Lean}
    \begin{itemize}
      \item $p^2<1$ hace positiva la base $1-p^2$.
      \item La potencia conserva positividad.
      \item Otro teorema prueba que $p^2=1$ produce una trayectoria plana.
    \end{itemize}

    \textbf{Resultado registrado}
    \begin{itemize}
      \item Nueve endpoints sin \texttt{sorry}.
      \item \texttt{check --fast}: PASS en WSL.
      \item Auditoría fuente--\texttt{Spec}: pendiente.
    \end{itemize}
  \end{columns}
\end{frame}

%================= CIERRE =============================================
\begin{frame}{12 \textperiodcentered{} Conclusiones}
  \begin{enumerate}
    \item El agente expande la frontera cuando añade una opción delegada con
      $T^D<T^S$ para lenguajes desconocidos.
    \item El estudio de GitHub encuentra la dinámica prevista, pero la selección
      en la adopción impide interpretar el evento como causalidad definitiva.
    \item La formalización cerró nueve resultados algebraicos y detectó una
      condición ausente: la estrictez de la Proposición 3 requiere $p^2<1$.
  \end{enumerate}

  \medskip
  \begin{center}
    \alert{La contribución verificable es el límite exacto del argumento, no una afirmación más amplia sobre lo que la IA siempre hace.}
  \end{center}

  \sourcefoot{Material analítico, derivación manuscrita y módulo Lean disponibles en el repositorio.}
\end{frame}

\end{document}
